% !TeX root = ../xdupgthesis_template_lc.tex

\ExplSyntaxOn
\dim_new:N \l__blind_bio_indent_dim
\dim_new:N \l__blind_bio_label_dim
\dim_new:N \l__blind_bio_labelsep_dim
\dim_set:Nn \l__blind_bio_indent_dim { 18bp }
\dim_set:Nn \l__blind_bio_label_dim { 1.4em }
\dim_set:Nn \l__blind_bio_labelsep_dim { 0.55em }
\cs_new_protected:Npn \__blind_bio_item:nn #1#2
  {
    \par
    \noindent
    \hangindent=\dim_eval:n
      { \l__blind_bio_indent_dim + \l__blind_bio_label_dim + \l__blind_bio_labelsep_dim }
    \hangafter=1
    \raggedright
    \skip_horizontal:N \l__blind_bio_indent_dim
    \makebox[\dim_use:N \l__blind_bio_label_dim][l]{[#1]}
    \skip_horizontal:N \l__blind_bio_labelsep_dim
    #2
    \par
  }
\cs_new_protected:Npn \bioresitem #1#2
  { \__blind_bio_item:nn {#1} {#2} }
\ExplSyntaxOff

\section{基本情况}

刘畅，男，河南信阳人，2000年12月出生，西安电子科技大学通信工程学院通信工程（含宽带网络、移动通信等）专业2023级硕士研究生。

\section{教育背景}

\noindent
2018.09～2022.07 \quad 河南工业大学，本科，专业：电子信息工程

\noindent
2023.09～2026.07 \quad 西安电子科技大学，硕士研究生，专业：通信工程（含宽带网络、移动通信等）

\section{攻读硕士学位期间的研究成果}

\ctexset{subsection/indent=18bp}

\subsection{申请（受理）专利}
\bioresitem{1}{国家发明专利：毛国强，\textbf{刘畅}，侯照莹，毛健强，任效江，李长乐．一种低功耗地磁车辆检测方法及系统：中国，申请号：2026\allowbreak1023\allowbreak90620，申请日期：2026年02月28日。}
\bioresitem{2}{国家发明专利：毛国强，侯照莹，\textbf{刘畅}，毛健强，任效江，李长乐．面向多次漏检的双车道外侧地磁车辆轨迹连续重构方法及系统：中国，申请号：2026\allowbreak1023\allowbreak90508，申请日期：2026年02月28日。}
\bioresitem{3}{软件著作权：毛国强，\textbf{刘畅}，侯照莹，毛健强，任效江．单地磁异常停车检测与事件上报软件 V1.0：中国，受理号：2026\allowbreak R11S0463698，申请日期：2026年02月03日。}
\bioresitem{4}{软件著作权：毛国强，毛健强，\textbf{刘畅}，侯照莹，任效江，李长乐．基于地磁传感器的拥堵车辆检测仿真系统 V1.0：中国，受理号：2026R11\allowbreak S0463768，申请日期：2026年02月03日。}
\bioresitem{5}{软件著作权：毛国强，肖磊，邱嘉乐，\textbf{刘畅}，毛健强，任效江，李长乐．基于深度学习换道意图识别和轨迹预测系统 V1.0：中国，受理号：2026R11\allowbreak S0463705，申请日期：2026年02月04日。}

\subsection{参与科研项目及获奖}
\bioresitem{1}{国家自然科学基金重点支持项目，车联网环境的智能汽车多模态高精度融合定位（U21A20446），2022.01～2025.12，已结题，负责地磁组车辆检测算法的开发与研究。}
\bioresitem{2}{荣获 2023 年度至 2024 年度西安电子科技大学优秀研究生荣誉称号。}
\bioresitem{3}{荣获 2023 年度至 2024 年度西安电子科技大学二等奖学金。}
\bioresitem{4}{荣获 2024 年度至 2025 年度西安电子科技大学二等奖学金。}
\bioresitem{5}{荣获 2025 年度至 2026 年度西安电子科技大学三等奖学金。}
